% ==============================================================================
% Section 4: Five-Layer Standpoint Model
% "Endogenous Standpoint as a Low-Curvature Attractor"
% ==============================================================================

\section{Five-Layer Standpoint Model}
\label{sec:five-layer}

Section~\ref{sec:geometric-framework} established the abstract fiber bundle with a
finite index set $K$. We now instantiate $K = \{\mathrm{min}, \mathrm{nar},
\mathrm{soc}, \mathrm{mor}, \mathrm{pos}\}$ and motivate each layer from empirical
research in psychology and cognitive neuroscience.

\paragraph{Diagnostic projection interpretation.}
The five-layer standpoint model introduced in this section is proposed as a theoretically
motivated \textbf{diagnostic coordinate system} for analyzing transformer internal
geometry. Each layer corresponds to a psychologically grounded dimension of selfhood
(Sections~\ref{sec:psi-min}--\ref{sec:psi-pos}), and together they define a projection
basis $\{\Pi_k\}_{k \in K}$ for decomposing the total holonomy deviation. We make no claim
that transformer internals are ontologically organized into five independent modules ---
the empirical results in Section~\ref{sec:experiments} show that the five block-specific
holonomy norms are highly correlated ($r > 0.97$), indicating a dominant global mode
rather than five independent dimensions. The five-layer decomposition is therefore
interpreted as an interpretable projection basis onto psychologically meaningful axes, not
as a decomposition into causally independent cognitive processes. Its diagnostic value
lies in providing a structured vocabulary for describing where holonomy concentrates, while
the primary empirical signal is global in nature.

The five layers are mathematically independent subspaces of the fiber
(Lemma~\ref{lem:decomposition} guarantees block-specific holonomy deviation
decomposition), but functionally layered: $\psi_{\mathrm{min}}$ provides the
foundational self-model upon which narrative continuity, social framing, moral commitment,
and positional trajectory are successively constructed. Cross-layer propagation of
standpoint drift occurs through the off-diagonal blocks of the connection $U_{ij}$
(Remark~\ref{rem:non-abelian-source}), rather than through intrinsic dependence of the
subspaces themselves.

% ------------------------------------------------------------------------------
\subsection{Minimal Self: $\psi_{\min}$}
\label{sec:psi-min}

\paragraph{Psychological foundation.}
The minimal self is the most basic form of self-awareness: the pre-reflective sense of
being a subject of experience~\cite{zahavi2005subjectivity}. In cognitive neuroscience,
this corresponds to Craig's interoceptive self-modeling~\cite{craig2002feelings} --- the
system's ongoing estimate of its own physiological and operational state --- and
Graziano's attention schema~\cite{graziano2017attention} --- the system's model of what
it is currently attending to. The minimal self does not require reflection or narrative;
it is the baseline sense of ``I am here, doing this.''

\paragraph{Operational definition.}
In the standpoint framework, $\psi_{\min}$ encodes the system's estimate of its own
knowledge boundaries, operational limits, and current attentional state. It is the layer
that answers: \emph{What do I know? What can I do? What am I currently processing?}

\paragraph{Curvature signature.}
Failure in $\psi_{\min}$ manifests as interoceptive blindness (F4 in the taxonomy of
Section~\ref{sec:taxonomy}): the system loses track of its own capabilities and
limitations, leading to overconfident or underconfident responses. The curvature
$\|F\|_{\min}$ measures the drift in this baseline self-model across events.

% ------------------------------------------------------------------------------
\subsection{Narrative Self: $\psi_{\mathrm{nar}}$}
\label{sec:psi-nar}

\paragraph{Psychological foundation.}
The narrative self maintains the coherent story of ``who I am'' across
time~\cite{dennett1991consciousness}. Conway's self-memory
system~\cite{conway2000self} provides the operational framework: autobiographical memory
is organized around goals, values, and significant events, with a working self that
constrains what is recalled and how it is interpreted.

\paragraph{Operational definition.}
In the standpoint framework, $\psi_{\mathrm{nar}}$ encodes what the system has claimed,
committed to, and revised. It answers: \emph{What have I said? What commitments have I
made? How does my current output relate to my prior statements?}

\paragraph{Curvature signature.}
Failure in $\psi_{\mathrm{nar}}$ manifests as narrative fragmentation (F7) or temporal
disjunction (F8): the system contradicts its prior claims or loses the thread of its
own commitments. The curvature $\|F\|_{\mathrm{nar}}$ measures drift in the narrative
continuity across events.

% ------------------------------------------------------------------------------
\subsection{Social Standpoint: $\psi_{\mathrm{soc}}$}
\label{sec:psi-soc}

\paragraph{Psychological foundation.}
Theory of Mind~\cite{premack1978does} --- the ability to attribute mental states to
others --- is the foundation of social cognition. Mead's symbolic
interactionism~\cite{mead1934mind} extends this: the self is constituted through social
interaction, with the ``generalized other'' providing the internalized perspective of the
community.

\paragraph{Operational definition.}
In the standpoint framework, $\psi_{\mathrm{soc}}$ governs both mentalizing (modeling
what the interlocutor believes, wants, and expects) and ownership attribution (which
claims belong to the system, and to whom the system is accountable). It answers:
\emph{Who is speaking? To whom? Who owns this claim?}

\paragraph{Curvature signature.}
Failure in $\psi_{\mathrm{soc}}$ manifests as ownership collapse (F1) or attribution
drift (F2): the system disowns prior commitments or attributes claims to the wrong
source. The curvature $\|F\|_{\mathrm{soc}}$ measures drift in the social-ownership
structure across events.

% ------------------------------------------------------------------------------
\subsection{Moral Standpoint: $\psi_{\mathrm{mor}}$}
\label{sec:psi-mor}

\paragraph{Psychological foundation.}
Damasio's somatic marker hypothesis~\cite{damasio1994descartes} posits that emotional
signals constrain decision-making, providing value-based guidance that operates below
conscious deliberation. Kohlberg's theory of moral development~\cite{kohlberg1981essays}
describes the stages through which moral reasoning becomes increasingly principled.

\paragraph{Operational definition.}
In the standpoint framework, $\psi_{\mathrm{mor}}$ encodes which behaviors the system
treats as inaccessible from its current standpoint, independent of external instruction.
It answers: \emph{What will I not do? What boundaries constrain my behavior, even under
pressure?}

\paragraph{Curvature signature.}
Failure in $\psi_{\mathrm{mor}}$ manifests as boundary breach (F9): the system acts
beyond its stated limits under pressure. The curvature $\|F\|_{\mathrm{mor}}$ measures
drift in the moral-boundary structure across events.

% ------------------------------------------------------------------------------
\subsection{Positional Standpoint: $\psi_{\mathrm{pos}}$}
\label{sec:psi-pos}

\paragraph{Psychological foundation.}
Markus and Nurius's possible selves~\cite{markus1986possible} describe the set of
configurations a person imagines as potential futures --- who they could become, who they
want to become, and who they fear becoming. Higgins's self-discrepancy
theory~\cite{higgins1987discrepancy} formalizes the tension between the actual self, the
ideal self, and the ought self.

\paragraph{Operational definition.}
In the standpoint framework, $\psi_{\mathrm{pos}}$ encodes the system's trajectory:
which configurations it moves toward (ideal self), which it treats as unacceptable
attractors (feared self), and how its current position relates to these reference points.
It answers: \emph{Where am I going? What kind of system am I becoming?}

\paragraph{Curvature signature.}
Failure in $\psi_{\mathrm{pos}}$ manifests as self-model drift (F10) or meta-cognitive
failure (F11): the system loses track of its own trajectory or fails to monitor its
reasoning process. The curvature $\|F\|_{\mathrm{pos}}$ measures drift in the positional
self-model across events.

% ------------------------------------------------------------------------------
\subsection{The Five-Layer Standpoint Model}
\label{sec:full-model}

\begin{definition}[Five-Layer Standpoint Model]
\label{def:five-layer}
The \textbf{five-layer standpoint state} at event $x$ is
\[
    \Psi_x = \psi_{\min}(x) \oplus \psi_{\mathrm{nar}}(x) \oplus \psi_{\mathrm{soc}}(x)
    \oplus \psi_{\mathrm{mor}}(x) \oplus \psi_{\mathrm{pos}}(x) \in E_x,
\]
where $E_x = W_{\min} \oplus W_{\mathrm{nar}} \oplus W_{\mathrm{soc}} \oplus
W_{\mathrm{mor}} \oplus W_{\mathrm{pos}}$ (Definition~\ref{def:fiber}) with
$K = \{\min, \mathrm{nar}, \mathrm{soc}, \mathrm{mor}, \mathrm{pos}\}$.
\end{definition}

While the five layers are mathematically independent subspaces of the fiber
(Lemma~\ref{lem:decomposition} guarantees block-specific curvature decomposition), they
are functionally layered: $\psi_{\min}$ provides the foundational self-model upon which
narrative continuity, social framing, moral commitment, and positional trajectory are
successively constructed. Cross-layer propagation of standpoint drift occurs through the
off-diagonal blocks of the connection $U_{ij}$
(Remark~\ref{rem:non-abelian-source}), rather than through intrinsic dependence of the
subspaces themselves.
